% 14_uso_ia.tex - Capítulo 14: Uso de IA Generativa
\chapter{Reporte de Uso de IA Generativa}
\label{ch:ia}

Este capítulo presenta de forma transparente el apoyo recibido de una herramienta de
inteligencia artificial en la redacción, organización, programación y revisión del trabajo.
La herramienta fue un apoyo para el equipo y no reemplazó el análisis, las decisiones,
la implementación ni la validación final del proyecto.

\section{Herramienta}
\label{sec:herramienta-ia}

\textbf{Herramienta de apoyo:} \textbf{Claude Code} (modelo \textbf{Claude Sonnet 5.5}),
utilizada por el equipo como apoyo para consultas técnicas, redacción, organización y
revisión del material. La herramienta no fue el desarrollador principal ni tomó decisiones
autónomas sobre el proyecto.

\section{Propósito y Áreas de Uso}
\label{sec:proposito-ia}

\begin{longtable}{p{3cm}p{3cm}p{8cm}}
\toprule
\textbf{Área} & \textbf{Qué hizo la IA} & \textbf{Qué hizo el equipo} \\
\midrule
\endfirsthead
\toprule
\textbf{Área} & \textbf{Qué hizo la IA} & \textbf{Qué hizo el equipo} \\
\midrule
\endhead
\midrule
\multicolumn{3}{r}{\textit{Continúa}} \\
\endfoot
\bottomrule
\endlastfoot
Contexto y proceso & Sugirió una organización de archivos y fases de trabajo & El equipo evaluó, adaptó y registró las decisiones en \texttt{decisions.md} \\
Decisiones & Presentó alternativas técnicas para discusión & El equipo eligió las reglas, el alcance y las convenciones del sistema \\
\code{domain/}, \code{scenario/} & Apoyó la redacción de borradores de código a partir de las reglas definidas & El equipo definió las reglas, revisó la implementación y exigió validaciones adicionales \\
BFS, DFS, UCS, A*, Beam, verifier & Apoyó la elaboración de código de referencia y pruebas & El equipo definió las convenciones y derivó los valores esperados antes de ejecutar \\
Heurísticas & Apoyó la exploración de alternativas y la redacción de ambas heurísticas & El equipo verificó admisibilidad, consistencia y dominancia h2$\ge$h1 \\
Utilidad, Minimax, $\alpha$-$\beta$ & Apoyó borradores de código y fórmulas para discusión & El equipo aprobó, corrigió y validó la utilidad y los algoritmos \\
CLI, runner, sweeps, export & Apoyó ejemplos y revisión de interfaces & El equipo decidió las opciones y los formatos CSV/JSON \\
GUI & Apoyó borradores de \code{gui/session.py} y \code{gui/app.py} & El equipo definió las funciones, revisó el comportamiento y validó la interfaz \\
Escenarios y experimentos & Apoyó la generación de mapas y la organización de resultados & El equipo escogió los casos y calculó los valores esperados \\
Documentación & Apoyó la redacción, corrección de estilo y organización de este informe y del manual & El equipo revisó, corrigió y asumió la revisión final del contenido \\
\bottomrule
\caption{Uso de IA por área}
\label{tab:uso-ia}
\end{longtable}

\section{Intervenciones Humanas y Correcciones}
\label{sec:intervenciones-humanas}

Las siguientes correcciones humanas se realizaron sobre los borradores de la IA:

\begin{enumerate}
\item \textbf{Signo del término rival (utilidad):} La decisión textual daba un único término
  negativo rival; la implementación hace que el signo dependa de quién porta el recurso
  (un rival cerca del \textit{portador de MAX} baja $U$; cerca del portador de MIN no suma).
\item \textbf{Rivales como obstáculos en búsqueda 1-agente:} El primer modelo trataba rivales
  como obstáculos fijos, lo que dejaba el Apéndice A sin solución (B1 sobre la meta). El
  equipo cambió a: \textbf{quitar rivales}, compañeros quietos como obstáculos (Sección 6.1).
\item \textbf{Inconsistencia de documentación:} Se revisaron los textos que describían el
  origen de las propuestas para distinguir claramente entre sugerencias de la herramienta
  y decisiones adoptadas por el equipo.
\item \textbf{Validación de casos 1--10:} Valores esperados derivados \textbf{a mano del mapa
  antes de ejecutar}; tests asercionan esos valores. La IA no inventó los resultados.
\item \textbf{Revisión humana:} El equipo revisó y validó los resultados, el código y la
  documentación generados con apoyo de herramientas de IA.
\end{enumerate}

La IA se utilizó como apoyo técnico y editorial; las decisiones de diseño, la
implementación, la verificación y la responsabilidad final del trabajo corresponden al
equipo.